\documentclass[10pt]{article}

\usepackage[utf8]{inputenc}

\usepackage[T1]{fontenc}

\usepackage{amsmath}

\usepackage{amsfonts}

\usepackage{amssymb}

\usepackage[version=4]{mhchem}

\usepackage{stmaryrd}

\usepackage{graphicx}

\usepackage[export]{adjustbox}

\graphicspath{ {./images/} }



\begin{document}

граммометр, прибавив соотношение: $I$ - тепловая энергия в калориях или механич. энергия в килограммометрах, где $I$ - размерная «физическая» постоянная, называемая механич. эквивалентом тепла.



Физич. законы, содержащие размерные постоянные, можно использовать для сокрацения числа первичных единиц измерения. Так, напр., постоянную скорости света можно считать абсолютной постоянной, т. е. безразмерной величиной, и таким путем выразить единицу измерения и размерность длины через единицу измерения и размерность времени; при рассмотрении гравитационной постоянной как абсолютного безразмерного числа можно выразить размерность и единицу измерения для массы через размерность и единицы измерения для длины $L$ и времени $T$.



В нек-рых вопросах явно проявляется тенденция к стандартизации путем введения единой универсальной системы единиц измерения. Однако привязывание универсальной системы единиц измерения к определенным физически фиксированным масштабам или к соответствующим физич. постоянным является искусственным. Наоборот, возможность использования произвольных единиц измерения и независимость рассматриваемых закономерностей от выбора систем единиц измерения могут служить источником ценных выводов.



Физич. закономерности, вообще говоря, не зависят от выбора системы единиц измерения. Это обстоятельство обусловливает особую структуру функций и функционалов, выражающих собой через посредство размерных величин физич. закономерности, независимые от систем единиц измерения. Эта специальная структура функциональных соотношений устанавливается $\pi-$ т е о р е м о й: всякое соотношение между размерными характеристиками, имеющее физич. смысл, представляет собой, по существу, соотношение между отвлеченными безразмерными комбинациями, к-рые можно составить из размерных определяемых и определяющих величин, среди к-рых должны учитываться и размерные физич. постоянные, имеющие существенное значение в рассматриваемых явления $x$.



Дополнительные данные об отсутствии нек-рых физич. постоянных в рассматриваемом соотвошении приводят к сокращению в этих соотношениях существенных аргументов. Напр., если в рассматриваемых тепловых И механич. явлениях нет преобразования - перехода тепловой энергии, измеряемой в калориях, в механическую, измеряемую в эргах, то постоянная механич. эквивалента тепла будет отсутствовать среди аргументов функции, описывающей соответствующую закономерность.



Для получения полезных выводов с помошцю Р. а. необходимо схематизировать проблему и, прежде всего, фиксировать общее моделирование явлений и свойств рассматриваемых объектов. Во многих случаях такая схематизация может быть связана с рядом рабочих гипотез. В рамках нек-рых моделей устанавливается система характеристик, к-рые связаны между собой физич. соотношением и к-рые согласно $\pi$-теоремө должны представляться как соотношения между безразмерными параметрами. Таким образом, нужно ввести систему определяющих параметров постоянных или переменных, вытекающую из постановки выделяемого класса задач и характеризующую, вообще говоря, полностью для данной среды каждую отдельно взятую задачу.



При выделении минимального числа определяющих параметров требуется учитывать условия симметрии и выбор выгодных систем координат. Независимые переменные (координаты точек пространсгва и времени) и физич. параметры типа коәфффициентов теплопроводности, вязкости, модулей упругости и т. П. необхо- димо включать в таблицу определяющих параметров. Такие постоянные, как ускорение силы тяжести или механич. эквивалент тепла, тоже должны фигурировать в перечне определяюпцих параметров, когда гравитация или переход тепловой энергии, измеряемой в калориях, в механическую, измеряемую в килограммометрах (илм в эргах), существенны.



Для выделенного класса задач необходимо включать в число определяющих параметров размерные или безразмерные характеристики заданных границ рассматриваемых тел и задаваемых функций, участвующие в формулировке начальных, краевых или каких-либо других условий.



Если рассматриваемая задача сформулирована как математическая, то можно выписать полную таблицу аргументов для соотношений вида



\[

a=f\left(a_{1}, a_{2}, \ldots, a_{n}\right),

\]



где $a$ - искомая величина, а $a_{1}, \ldots, a_{n}$ - определяющие параметры. В зависимости от постановки задачи число $n$ может быть конечным или бесконечным. Обычно при соответствующем фиксировании класса рассматриваемых задач можно огранитиваться случаями, когда $n$ конечно и, вообще говоря, невелико. Параметры $a_{1}, \ldots, a_{n}$ составляются из задаваемых величин, входяших в основные уравнения, граничные условия и начальные данные. Определяющие параметры -. это все исходные данные, к-рые надо знать предварительно по смыслу постановки математич. задачи для вычисления искомой функции различными путями, в том числе и при расчетах на специальных машинах или с помощью аналоговых систем.



Определяющие параметры при получении нужных ответов с помощью экспериментов - это велічины, характеризующие каждый отдельный опыт, величины. к-рые необходимы и достаточны для повторения и сравнения различных әкспериментов.



Можно выписать систему определяющих параметров и в тех случаях, когда детальные свойства модели и системы уравнений, описывающих рассматриваемые явления, вообще говоря, неизвестны. достаточно опереться на предварительные данные или гипотезы о виде задаваемых функций и о постоянных, к-рые входят или могут входить в определение модели, и на др. условия, выделяющие конкретные решения задачи.



Выводы Р. а. получаются на основании $\pi$-теоремы о существовании соотношения вида (3), к-рое выполняется для определяемой величины и определяюших величин $a_{1}, \ldots, a_{n}$ в любой системе единиц измерения.



Система аргументов в соотношении (3) должна в смысле теории размерностей обладать полнотой. Это значит, что если величина $a$ отлична от нуля или бесконечности, то существуют числа $\alpha_{i}, \ldots, \alpha_{n}$ такие, что размерности величины $a$ и комбинаций $a_{1}^{\alpha_{1}} a_{2}^{\alpha_{2}} \ldots a_{n}^{\alpha_{n}}$



\includegraphics[max width=\textwidth, center]{2023_07_08_bd46520950157045bb54g-1}

шение (3) можно переписать в виде



\[

\pi=\frac{a}{a_{1}^{\alpha_{1}} a_{2}^{\alpha_{2}} \ldots a_{n}^{\alpha_{n}}}=g\left(\pi_{1}, \pi_{2}, \ldots, \pi_{n-k}\right),

\]



где $\pi_{1}, \ldots, \pi_{n-k}$ - степенные безразмерные одночлены' (аналогичные $\pi=3,14$. .), образованные из $a_{1}, \ldots, a_{n}$, причем $k \leqslant n$. Число $k$ равно числу параметров среди $a_{1}, \ldots, a_{n}$ с независимыми размерностями .



После явной записи физич. закономерностей в виде (4) по существу заканчиваются возможности применений P. а. Соответствующие результаты, получевные таким путем, носят ограниченный характер. Если видоизменить любым способом уравнения движения без внесения в них каких-либо новых параметров, то физич. закономерности (4) могут сильно измениться,





\end{document}